\section{Problem set}
\subsection{Discussion 3 and related Lecture 10 practice}
\begin{prob}
\textbf{Q1: Swipe.} Print the digits of a positive integer $n$ backward and then forward, one per line, with the leftmost digit printed once. No \py{while}, \py{for}, or \py{str}. For 2837 the output is 7,3,8,2,8,3,7. The supplied base is \py{if n<10: print(n)}; fill its else branch.\par
\hint{Use the before/call/after arrangement from cascade, printing just the last digit.}
\end{prob}
\begin{prob}
\textbf{Q2: Skip Factorial.} For positive integer $n$, return $n(n-2)(n-4)\cdots$ down to 1 or 2. Fill the base condition, base return and recursive return. Source checks:5 gives 15;8 gives 384.\par
\hint{One base condition can include both possible final factors.}
\end{prob}
\begin{prob}
\textbf{Q3: Is Prime.} For integer $n>1$, implement a recursive primality test with no while/for. Source tests:2 true,16 false,521 true. Use a helper with as few arguments as possible, and give its docstring a precise behavioral description. The suggested template is:
\begin{lstlisting}
def is_prime(n):
    def f(i):
        if i == n:
            return __________
        elif ________________:
            return __________
        else:
            return f(________)
    return f(2)
\end{lstlisting}
\hint{Retain n in the outer frame; the helper searches candidate divisors from i through n-1.}
\end{prob}
\begin{prob}
\textbf{Q4: Recursive Hailstone.} For positive integer $n$, print every term through 1 and return the number of elements, including the first and last. For 10, output 10,5,16,8,4,2,1 and return 7; at 1 print 1 and return 1. Preserve the supplied split into hailstone/even/odd:
\begin{lstlisting}
def hailstone(n):
    print(n)
    if n % 2 == 0:
        return even(n)
    else:
        return odd(n)

def even(n):
    return __________________

def odd(n):
    # Complete this helper.
    _________________________
\end{lstlisting}
\hint{An even input always continues; an odd input has a special base at 1. Count the current term in each helper's returned value.}
\end{prob}
\begin{prob}
\textbf{Q5: Sevens.} Implement \py{sevens(n,k)} for positive integers $n,k$, returning which of the $k$ numbered players says $n$ under the direction rule developed above. The source requires helper \py{f(i,who,direction)}, with base \py{if i==n:return who} and initial call \py{f(1,1,1)}; supplied has\_seven removes digits recursively. Source inputs $n=2,6,7,8,9,18$, each with $k=5$, give 2,1,2,1,5,2.\par
\hint{Switch the direction at the current spoken number, advance/wrap who, then recurse on i+1.}
\end{prob}
\begin{prob}
\textbf{Spring 2024 MT1 Q4(e), reproduced in Lecture 10 slide 6: Streak.} A dice integer has every digit in 1--6. For positive $n$, return whether it is a dice integer with every digit the same. Fill the two source blanks:
\begin{lstlisting}
def streak(n):
    return (n >= 1 and n <= 6) or (
        n > 9 and _______________ and streak(________))
\end{lstlisting}
Source true cases 22222,4,33333333333333; false 22322,99999,505,707,7070.\par
\hint{A one-digit valid die is a base; a longer number needs equal last two digits and a valid shorter streak.}
\end{prob}
\begin{prob}
\textbf{Q6 (optional): Karel the Robot.} Karel starts at the bottom-left corner of an unknown $n\times n$ square. Use the east-facing start of the \href{https://compedu.stanford.edu/karel-reader/docs/python/en/ide.html}{linked default UI}. Four available functions: \py{move()} advances one square or errors at a wall; \py{turn\_left()} turns 90 degrees left; \py{front\_is\_blocked()} and \py{front\_is\_clear()} test the square ahead. End in the middle of the bottom row, in any orientation: for 7, column 4; for even $n$, either column $n//2$ or $n//2+1$.

Only if/if-else statements and function calls may appear in main's body. No assignment, nested def, lambda, while or for. You may call main recursively. The discussion prose omits the initial heading; the linked UI's bottom-left/east default supplies the orientation used here.\par
\hint{Advance two squares before the recursive call, then return one square while coming back. Specify the direction in which each smaller call finishes.}
\end{prob}

\subsection{Homework 3 --- due September 24 at 11:59p.m.}
\begin{prob}
\textbf{Q1: Num Eights.} For positive integer $n$, return its number of 8 digits using recursion. No assignment statements of any kind or loops, including in helpers. Source checks:3 $\mapsto$0,8 $\mapsto$1,88888888 $\mapsto$8,2638 $\mapsto$1,86380 $\mapsto$2,12345 $\mapsto$0,8782089 $\mapsto$3.\par
\hint{Separate the last digit's contribution from the count on the remaining prefix.}
\end{prob}
\begin{prob}
\textbf{Q2: Digit Distance.} Return the sum of absolute differences between consecutive digits of positive integer $n$. Use recursion and no loops. Source illustrations:61 gives 5;71253 gives $6+1+3+2=12$;6 gives 0. Starter tests:3 $\mapsto$0,777 $\mapsto$0,314 $\mapsto$5,31415926535 $\mapsto$32,3464660003 $\mapsto$16.\par
\hint{Removing the last digit removes exactly one adjacent pair.}
\end{prob}
\begin{prob}
\textbf{Q3: Interleaved Sum.} For positive integer $n$, compute $f_{\mathrm{odd}}(1)+f_{\mathrm{even}}(2)+f_{\mathrm{odd}}(3)+\cdots$ through $n$. Do not directly test parity. No loops, \py{\%}, \py{//}, multiplication, or bitwise \py{\& | \textasciicircum} in the function or its helpers; the public tests ban these syntax nodes individually. With identity, square and triple, source cases are $(5,\mathrm{identity},\mathrm{square})=29$, $(5,\mathrm{square},\mathrm{identity})=41$, $(4,\mathrm{triple},\mathrm{square})=32$, $(4,\mathrm{square},\mathrm{triple})=28$.\par
\hint{The first index is known to be odd; alternate two helpers whose names determine which function to apply.}
\end{prob}
\begin{prob}
\textbf{Q4: Count Coins.} Count unordered ways to make positive integer total using 1,5,10,25-cent coins. Use recursion, no loops, and the supplied next\_smaller\_coin:25 $\to$10 $\to$5 $\to$1 $\to$ None. At 15 there are six ways: fifteen 1s; ten 1s+one 5; five 1s+two 5s; five 1s+one 10; three 5s; one 5+one 10. Starter totals 15,10,20,45,100,200 give 6,4,9,39,242,1463.\par
\hint{Use the include/exclude recurrence at the current denomination; test for None before subtracting a coin value.}
\end{prob}
\begin{prob}
\textbf{Q5: Maximum Subsequence.} For a nonnegative integer $n$ and nonnegative integer limit $t$, return the largest numeric value obtainable by keeping at most $t$ digits in their original order. Use recursion, no loops. A subsequence need not be contiguous. Source checks: $(2012,2)\mapsto22$, $(20125,3)\mapsto225$, $(20125,5)$ and $(20125,6)$ both give 20125; $(12345,3)\mapsto345$, $(12345,0)\mapsto0$, $(12345,1)\mapsto5$. Numeric zero represents the empty choice as well as all-zero choices.\par
\hint{Either append the last digit to the best shorter subsequence with t-1 slots, or skip it and keep t slots.}
\end{prob}
\begin{prob}
\textbf{Q6 (just for fun, outside required course scope): Towers of Hanoi.} Move a stack of $n\ge1$ disks from rod start to a different rod end among 1,2,3. Initially all disks are on start; smaller disks are above larger ones and the other rods are empty. Move one top disk at a time, never onto a smaller disk. Use the given print\_move function to print each move, and retain the supplied valid-rod assertion. The starter also describes subproblems where larger disks remain under the stack being moved.\par
\hint{Move the top n-1 disks to the spare rod, move the largest disk, then move the n-1 disks to the destination.}
\end{prob}
\begin{prob}
\textbf{Q7 (just for fun, outside required course scope): Anonymous Factorial.} Return one expression computing a factorial function, using only calls, conditional expressions and lambdas. No assignment/def statements in the returned expression, and do not refer to make\_anonymous\_factorial. Supplied imports are sub and mul from operator. Source check: the returned function at 5 gives 120.\par
\hint{A parameter can give the recursive function a local name; pass that function itself as an argument.}
\end{prob}
\par\noindent\begin{minipage}{\linewidth}
Check with \py{python3 -m pytest -k num\_eights}, replacing the final name for each other function. Prepare the Provenance submission bundle in VS Code and upload its zip to Gradescope. The official assignment also requires a checkoff interview; the just-for-fun puzzles are explicitly outside required course scope. Hog continues from Week 2 and is due September 17; no new lab is assigned this week.
\end{minipage}\par

\subsection{Assigned real exam practice}
\begin{prob}
\textbf{Fall 2017 MT1 Q4(a): Digital, p5 (linked by HW3).} Implement collapse for a nonnegative integer $n$, removing any digit equal to the digit immediately to its right, so no adjacent digits agree. Source checks:1234 $\mapsto$1234;12234441 $\mapsto$12341;0 $\mapsto$0;3 $\mapsto$3;11200000013333 $\mapsto$12013. Fill the original template:
\begin{lstlisting}
def collapse(n):
    left, last = n // 10, n % 10
    if ______________________:
        return last
    elif ______ == __________:
        return collapse(_____)
    else:
        return collapse(_____) * 10 + _____
\end{lstlisting}
\hint{Compare the last two digits; either discard a duplicate or append a different last digit after collapsing the prefix.}
\end{prob}
\begin{prob}
\textbf{Fall 2019 Final Q6(b): Palindromes, p6 (linked by HW3).} Implement contains(a,b) to test whether the digits of nonnegative integer $a$ appear in order among the digits of nonnegative integer $b$. No \py{str, repr, list, tuple, [, ]}. Source checks: \py{contains(357,12345678)} true, \py{contains(753,12345678)} false, \py{contains(357,37)} false. Follow the official solution's convention that $a=0$ denotes no remaining digits, including an initial 0; it is not a requirement to find the written digit 0.
\begin{lstlisting}
def contains(a, b):
    if a == b:
        return True
    if ______ > ______:
        return False
    if ______ == ______:
        return contains(______, ______)
    else:
        return contains(______, ______)
\end{lstlisting}
\hint{Match the requested last digit to the rightmost available matching digit of b, or discard a nonmatching last digit of b.}
\end{prob}

\textbf{Attribution.} Official Fall 2026 CS 61A lectures/assignments; Composing Programs 3ed section 1.7, CC BY-NC-SA 4.0. Original assigned exams: \href{https://cs61a.org/resources/fa17/mt1/61a-fa17-mt1.pdf}{Fall 2017 Midterm 1} and \href{https://cs61a.org/resources/fa19/final/61a-fa19-final.pdf}{Fall 2019 Final}. Other exam reproductions are identified above.
